% Source: https://artofproblemsolving.com/community/c3h2416904_conjecture_that_probably_exists_and_is_provendisproven
% Stable identity: aops-topic-2416904-q01
\textbf{(a)} Take a circle and one of its diameters. They cut the plane into three regions: the upper half-disc $1$, the lower half-disc $2$ and the outside $3$. Regions $1$ and $2$ share the diameter, and region $3$ shares an arc of the circle with each of them. So the three regions are pairwise adjacent, and two colours cannot work (two of the three regions would get the same colour).

\begin{asy}
usepackage("amsmath");
size(3cm);
defaultpen(fontsize(10pt));
fill(arc((0,0),1,0,180)--cycle,rgb(0.6,0.6,0.6)); fill(arc((0,0),1,180,360)--cycle,rgb(0.85,0.85,1));
draw(circle((0,0),1),linewidth(0.9)); draw((-1,0)--(1,0),linewidth(0.9));
label("$1$",(0,0.45)); label("$2$",(0,-0.45)); label("$3$",(1.15,0.9));
draw(box((-1.5,-1.3),(1.5,1.3)),gray+dotted);
\end{asy}

\textbf{(b)} Let the lines be $\ell_1,\dots,\ell_k$, and for each line call one of its two sides positive. For a point $X$ on none of the lines, let $c(X)$ be the number of lines for which $X$ is on the positive side. A region does not meet any line and is connected, so it lies on one side of each line, and $c$ is constant on it. Colour a region black if $c$ is even and white if $c$ is odd.

Suppose two regions share a boundary segment. The segment lies on one of the lines, say $\ell_i$ (it cannot lie on two different lines, which share at most one point). Choose a point $Y$ inside the segment that is not on any other line; this is possible because the other lines meet $\ell_i$ in finitely many points. Points close to $Y$ on the two sides of $\ell_i$ lie in the two regions and are on the same side of every line except $\ell_i$. So the values of $c$ in the two regions differ by exactly $1$, and the regions have different colours. $\blacksquare$

\begin{asy}
usepackage("amsmath");
size(3.5cm);
defaultpen(fontsize(10pt));
real L=10; pair[] P={(0,0.3),(0,-0.2),(0.4,0)}; pair[] D={dir(20),dir(125),dir(80)};
path hp(pair p, pair d){ pair n=rotate(90)*d; return (p-L*d)--(p+L*d)--(p+L*d+L*n)--(p-L*d+L*n)--cycle; }
fill(hp(P[0],D[0])^^hp(P[1],D[1])^^hp(P[2],D[2]),evenodd+rgb(0.6,0.6,0.6));
for(int i=0;i<3;++i) draw((P[i]-L*D[i])--(P[i]+L*D[i]),linewidth(0.9));
clip(box((-1.5,-1.3),(1.5,1.3)));
draw(box((-1.5,-1.3),(1.5,1.3)),gray+dotted);
\end{asy}

(The figure shows such a colouring for three lines: the shaded regions are on the positive side of an odd number of lines.)
